\section{Determination of the constant and covariance of its realizations}
\label{act21:centralidad-constante}

The holographic arithmetic--geometric coupling constant combines
structural determination with a capacity for reconstruction. The former
fixes the ordered register compatible with the incidences and chronology
of the generated state. The latter establishes how that register is
recovered from its representations and how operations acting on it are
identified. Their composition gives a precise meaning to its coupling
function.

Determination begins with the regions produced by APP--TRIT--TPK. Their
six-trit bands retain the position of every symbol, charge, carry and
depth. The integer clock compares the ternary and decimal scales; its
first zero increment, followed by the nonadic return, determines the
phase--charge reading of the descriptor. The terminal repertoire and its
incidences retain that provenance. Orderings of this repertoire produce
the finite domain on which the joint selection of
Section~\ref{act21:chap:despliegue-selector-k} acts.

The exceptional condition restricts that domain to registers whose
support face agrees with the face extracted from the regional words.
The temporal condition compares, within a single ordering, the
comparative and affine readings of the two terminal segments. Thus a
coordinate position and the relation between two instants enter the same
decision. The exact result is
\[
 19\,446\longrightarrow79\longrightarrow1.
\]
The unique surviving register determines the integer
\[
 N_K=234543140729659824621058914794146601,
 \qquad
 \kappa_{\rm ag}=\frac{N_K}{10^{36}-1}.
\]
The rational expression preserves the order of the twelve triads in the
specified phase chart. Uniqueness concerns the constructed terminal
domain, its two explicit conditions and its orientation data. This fixes
what is selected, through which relations, and within which set of
alternatives.

The census of the three regional endings serves a different function.
Its \(196\cdot197\cdot196\) elements are triples of six-trit words
compatible with the declared cylinders. The Gram profile, weights,
distances and oriented margins determine the visible ending. The census
of \(19\,446\) elements examines ordered twelve-coordinate registers.
Both receive information from the same enriched history; their types and
filters make it possible to follow their relation while distinguishing
a visible triple from the complete register.

Once \(K\) has been selected, the Hadamard transformation and its readers
establish equivalences between representations of that same register.
Selection determines a candidate among alternatives; inversion recovers
it from a representation. The bilateral law then provides the
conservation identity for the declared isometries, and the archive retains
the information needed to invert successive readings. Each result passes
to the next an object with specified domain, orientation and normalization.

This hierarchy also specifies the role of \(K\) in the conservation law.
The bilateral identity is proved for two isometries with a common domain
and codomain. In the realizations developed here, the selected register
fixes the arithmetic and incidence readings, its centred direction and
the reconstruction operations composed with that identity. Generality
of the balance and determination of the register are complementary
properties: the former describes an operator law; the latter identifies
the structure on which the present realizations are calculated.

\pagebreak % REV03: balance the three complete concluding paragraphs.
The closure of the fine-structure constant uses the already-generated
regional coordinates and the selected register, preserving the carries
of the dodecaphase operation. Its analytic representation develops the
same precoordinate through the complete recurrence. Through the other
projection, the preserved incidences are transported to the exceptional
realizations. Their common provenance relates arithmetic value, incidence
and memory while preserving the distinct quantities read by their
respective operators.

The subsequent geometric realizations extend this function. The centred
direction of \(K\) determines the determinant-one diagonal flow,
preserves Plücker signs and the covolume of the specified lattices; the
archive reconstructs that direction and controls its readout errors.
The theta--Mellin transformation preserves the polar structure on the
proved domain, while the entire difference records spectral variation.
Through the operations already established, temporal selection is thus
linked to a structure admitting evaluation, transport and recovery.

The combined contribution is therefore organized into a definite proof
sequence: generation of the state, selection of the constant, equivalence
of its representations, operator conservation and stable reconstruction.
This sequence fixes the position of each theorem and keeps visible both
the conditions it receives and the relations it transmits.
